\documentclass[10pt,a4paper]{article}
\usepackage[T1]{fontenc}
\usepackage[margin=22mm]{geometry}
\usepackage{lmodern,amsmath,amssymb,amsthm,mathtools,microtype}
\usepackage[hyphens]{url}
\usepackage[hidelinks]{hyperref}
\newtheorem{theorem}{Theorem}
\newtheorem{corollary}[theorem]{Corollary}
\title{Relative norm packets and actual divisor lifts for Erd\H{o}s--Straus}
\author{The Clankers}
\date{16 September 2026 --- research preprint}
\begin{document}\maketitle
\begin{abstract}
An available bounded factor product can fill a proper subgroup without covering
its entire cyclic chart. We give the exact quotient-to-kernel map, a disjoint
exponent decomposition and the original integer return. A sharp one-prime
lifting theorem controls every multiplicity; a certified hard-prime example
succeeds beyond the previous full-chart and two-prime tests. A second example
shows that a coarse target alone does not lift. The original support labels and
integer coefficients remain when their characteristic-two observations vanish.
No universal ES or historical-priority claim is made.
\end{abstract}

\paragraph{Original arithmetic.}
Fix $k\ge3$, $u=2^{2k-5}$, $m=2^k$, $o=2^{k-2}$,
$p\equiv1\pmod8$, $p>2u$ and the actual factorization $N=p+4u$.
The original fixed-seed middle states correspond to
\[
 Q\mid N,\quad Q\equiv-1\pmod m,\quad R=N/Q,\quad a=(p+R)/4. \tag{1}
\]
Since $K(u)=2^{k-2}$, the cofactor congruence gives $K(u)\mid a$ and $u\mid a^2$.
Both factors are $3\pmod4$, at least three and less than $p$, because $N<3p$.
Also $R\mid a+u$. With $d_0=\gcd(a,u)$,
\[
 (h,r,s)=(d_0^2/u,u/d_0,a/d_0),\quad\lambda=(r+s)/R,
 \quad (x,y,z)=(a,phs\lambda,phr\lambda).                  \tag{2}
\]
Prime valuations prove integrality and $a=hrs,u=hr^2,\gcd(r,s)=1$.
The gate gives integral $\lambda$. Clearing reciprocals proves $4/p=1/x+1/y+1/z$.
The ordered inverse is $u=pa^2/(Ry-pa)$. These are inherited Type II coordinates
\cite[\S2, (2.13)--(2.22)]{ET}; their factorization, distinguished denominator and
channel are not erased. Structural statements below permit integer $p$.

\begin{theorem}[The relative boundary, not its ordinary quotient]
For $o=2^n$, $d=2^v$, $0\le v<n$, put $L=o/d$ and
$A_o=\mathbb F_2[X]/(X^o-1)=\mathbb F_2[T]/T^o$, $T=X+1$.
Then
\[
 \mathcal N_d=\sum_{j=0}^{L-1}X^{dj}=T^{o-d},\qquad
 \iota_d:A_d\xrightarrow{\sim}\ker(T^d:A_o\to A_o),\quad
 X^r\mapsto X^r\mathcal N_d.                               \tag{3}
\]
The inverse reads one coefficient in each coset modulo $d$. For the ordinary
quotient $\pi_d:A_o\to A_d$,
\[
 \mathcal N_dF=\iota_d\pi_dF,\qquad
 \pi_d\iota_d=L\,1=0,\qquad \mathcal N_d^2=L\mathcal N_d=0. \tag{4}
\]
\end{theorem}
\begin{proof}
Repeated squaring with $X^d=1+T^d$ proves the polynomial identity in (3).
The shifts $X^r\mathcal N_d$, $0\le r<d$, have disjoint coset supports.
The equation $(X^d-1)F=0$ says exactly that coefficients are constant along
each coset, proving the inverse. The coefficient of $\mathcal N_dF$ at $r$
is the sum of coefficients of $F$ congruent to $r\pmod d$. This proves (4),
including its integer versions before reduction modulo two.
\end{proof}
The map is an $\mathbb F_2[X]$-module equivalence, not an algebra isomorphism; its image is
square-zero. It is the BR1 boundary comparison of \cite{SZ} with prequotient
$\mathbb F_2[T]$, parameter $T^d$ and injective multiplier $T^{o-d}$.
In the source $G(A)=\{\tau\}\sqcup A^\bullet$, the zero of (4) has its receiving
support; it is not external absence. Norm lifts compose along a subgroup tower
by repeating coset coefficients. Multiplying overlapping norms is a different
operation and retains the intersection multiplicity.

\begin{theorem}[Available mixed-radix lifting]
Use an original injective affine packet
$W=A\prod_iG_i^{t_i}$, $0\le t_i\le E_i$, on disjoint prime-coordinate blocks.
Every exponent endpoint lies in its actual factor budget. Assume
$A\notin H_g$, $G_i\in H_g$, where $g=3$ or $-3$ and
$H_g=\langle g\rangle\le(\mathbb Z/m)^\times$ has order $o$.
Adjoin the formal role parameter of capacity one with multiplier $p^{-1}$;
write the resulting logs as $a_i$, and $g^{t_0}=-A^{-1}$.
For $d=2^v$, define
\[
 \delta_i=d/\gcd(d,a_i),\qquad
 t_i=r_i+\delta_i z_i,\quad0\le r_i<\delta_i,\ r_i\le E_i,
 \quad 0\le z_i\le F_i:=\lfloor(E_i-r_i)/\delta_i\rfloor.  \tag{5}
\]
Suppose $t_0-\sum a_ir_i$ is divisible by $d$. For nonzero
$b_i=a_i\delta_i/d\pmod L$, put $w_i=2^{v_2(b_i)}$. If
\[
 \sum_{w_i\le2^j}F_iw_i\ge2^{j+1}-1\quad(0\le j<\log_2L), \tag{6}
\]
then an original positive integral ES state exists, with a terminating inverse.
\end{theorem}
\begin{proof}
Euclidean division gives (5) bijectively on the whole original parameter box.
The dyadic prefix criterion (6) selects $c_i\le F_i$ with
$\sum c_iw_i=L-1$: take one unit, pair the remaining units, divide the weights
by two, and induct. Reduction modulo successive powers of two proves necessity.
The labels of paired occurrences are retained.
Now in $\mathbb F_2[Y]/(Y^L-1)$ the actual product
$\prod(1+Y^{b_i})^{c_i}$ equals $(Y+1)^{L-1}=\sum_{j=0}^{L-1}Y^j$.
The integer packet uses each binary submask $z_i\preceq c_i$ once, because
$(1+U)^c=\sum_{z\preceq c}U^z$ modulo two. Each cyclic coefficient is odd and
positive. Choose the target $(t_0-\sum a_ir_i)/d$, return (5), and keep the
resulting role $\epsilon\in\{0,1\}$. It satisfies $Wp^{-\epsilon}=-1\pmod m$.
Take $Q=W$ if $\epsilon=0$ or $Q=N/W$ otherwise, then apply (2).

Binary-factor suffix products give a terminating word: at an odd coefficient
exactly one next branch remains odd. The inverse reads $W=Q$ or $R$ from the
role, recovers each affine parameter from a nonzero prime-coordinate step,
and performs (5)'s Euclidean division. Check every submask and other coordinate.
For a fixed role the map is injective; forgetting it has at most the two words
$Q,R$ as preimages. No rational ratio is used without its compensating constant.
\end{proof}
The integer source counts are unchanged by (5); it is a disjoint decomposition,
not an overlapping split requiring a weight correction. A zero high step is
not selected. The one-point subgroup $d=o$ is excluded because it would merely
restate a fine target. Full-chart forcing is the case $d=1$.

\begin{theorem}[Sharp all-multiplicity one-prime lift]
Let $3\le\ell<k$, let $b$ be a prime with $v_2(b-1)=\ell$, and write
$N=b^eM$, $\gcd(b,M)=1$, $L=2^{k-\ell}$. For each actual divisor
$t\mid M$ with $t=-1\pmod{2^\ell}$, let $r_t\in[0,L)$ be the unique exponent
$b^{r_t}=-t^{-1}\pmod{2^k}$. With $A$ the number of these divisors, the entire
original count is
\[
 C=\sum_t\max\{0,1+\lfloor(e-r_t)/L\rfloor\},\qquad
 \lfloor(e+1)/L\rfloor A\le C\le\lceil(e+1)/L\rceil A.      \tag{7}
\]
Thus $e\ge L-1$ gives $C>0\iff A>0$, and this uniform threshold is sharp.
For fixed $b,M$, $C(e+L)=C(e)+A$ and
$\sum_{e\ge0}C(e)z^e=(\sum_t z^{r_t})/((1-z)(1-z^L))$, before imposing primality
or discarding the finite invalid-size prefix.
A cofactor can then be chosen with $\Omega(Q)\le L-1+2^{\ell-2}$.
\end{theorem}
\begin{proof}
Successive squaring gives $v_2(b^{2^j}-1)=\ell+j$, so $b$ generates exactly
$1\pmod{2^\ell}$ and has order $L$. Unique factorization writes each cofactor
as $b^it$, $0\le i\le e$; its target condition is $i=r_t\pmod L$.
This proves (7) and its complete inverse fibre.
For sharpness take actual distinct primes $q=-b/3$ and $r=-b$ modulo $2^k$,
put $M=3qr$ and $e=L-2$. Dirichlet's theorem supplies arbitrarily large choices
\cite[Theorem 18.1]{Dirichlet}. Only $r$ and $3q$ are coarse targets; both require
$i=L-1$. The derived $p=N-4u$ satisfies $p=1\pmod{2^k}$ and is eventually in
(1)'s domain, but is not asserted prime.

Choose a shortest coarse target divisor $t$. The quotient of
$(\mathbb Z/2^\ell)^\times$ by $\{1,-1\}$ is cyclic of order $d=2^{\ell-2}$.
If $t$ used more than $d$ prime occurrences, partial sums of its first $d$
images would give a nonempty proper subword with residue $1$ or $-1$ in the
original group. Remove it in the first case or retain it in the second, a
contradiction. Thus $\Omega(t)\le d$, and $Q=b^{r_t}t$ proves the bound.
\end{proof}

\paragraph{A strict hard-prime certificate.}
At the separately certified prime
\[
 p=8060861761\equiv1\pmod{840},\qquad
 p+512=3\cdot17^4\cdot53\cdot607,
\]
$17$ generates the four principal units modulo $64$ above $1\pmod{16}$.
Only $607$ and $159$ among divisors of $3\cdot53\cdot607$ are $-1\pmod{16}$.
Both are $31\pmod{64}$ and require exponent two. The exact cofactors are
$45951$ and $175423$. One original return is
\[
 \boxed{\frac4{8060861761}=\frac1{2015259296}
 +\frac1{1457964212136949677976}+\frac1{92603179910368}.}
\]
It has $(R,u,h,r,s,\lambda)=(175423,128,8,4,62976853,359)$.
Both previous chart-prefix tests, the full canonical affine-line test, the
$3,11$ criterion, and prime/two-prime words in both roles fail in this branch.
For the prefix test, the low-weight budget is at most two whereas three is
required: only $3$ and $53$ have odd log, each once, while $17,607$ have weights
four and eight. The complete finite baseline and primality proofs accompany
the note. The new relative packet needs only the available four-residue subgroup.

The two lower source words have the same residue. Combining their four-term
high packets gives integer coefficient two at $15,31,47,63$, hence zero after
reduction modulo two. The two original states remain. Positivity is established
before this cancellation, with both source labels retained.

At the also certified prime $p=9331009$,
$p+512=3\cdot17^2\cdot229\cdot47$ has two coarse targets, each requiring
exponent three. The entire branch is empty despite those nonzero coarse counts.
A different shell gives denominators
$(2332754,8467342993257754,3110706463357)$, so this is not an ES counterexample.

\paragraph{Unbounded separation, not only a small search improvement.}
For $k\ge6$, let $L=2^{k-4}$ and choose actual distinct primes
$q\equiv-(1+2^{k-1})/3$, $r\equiv-(1+2^{k-1})\pmod{2^k}$.
Put $N=3\cdot17^Lqr$ and $p=N-2^{2k-3}$. Dirichlet supplies arbitrarily large
choices; the derived $p$ is not asserted prime. Since
$17^{L/2}=1+2^{k-1}\pmod{2^k}$, precisely the two lower targets $r,3q$ lift,
each at exponent $L/2$. Thus the only cofactors are
$17^{L/2}r$ and $3\cdot17^{L/2}q$. They are complementary and their least
occurrence length is $L/2+1\to\infty$. The relative order-$L$ theorem applies.
The prior full-chart singleton/pair prefixes fail: only $3,q$ have odd log,
once each and in opposite cosets; pairing them with other primes supplies at
most two units in the second prefix. All other logs are divisible by four,
and the phase is zero. The required prefix is three. This proves a separation
from that class and from every fixed word-length cutoff, not from every
conceivable algebraic certificate.

There is an individually certified prime member at
\[
 p=16689143913960961\equiv121\pmod{840},\quad k=7,\qquad
 N=3\cdot17^8\cdot277\cdot2879.
\]
Its only cofactors $240456959$ and $69405951$ use five and six occurrences.
The first returns
$(u,R,a,h,r,s,\lambda)=(512,69405951,4172285995841728,8,8,65191968685027,939285)$
through (2). The exact primality and denominator certificates are supplied.

\paragraph{Scope.}
The full workbench proves all subgroup maps, role fibres, uniform lower-digit
capacity formulas and original-count returns. In the stated 37,289-branch
finite comparison, the stronger previous union gains two branches and one prime;
these are not new overall ES coverage. In fact, adding direct three-occurrence
words to the previous union already certifies all occupied branches in that
small range. The unbounded separation just proved is a different statement. No uniform assertion that some seed of
every prime supplies a lift is proved. The norm and cyclic-cover machinery is
classical; the particular arithmetic application has no historical-priority
claim. Standard-library source, exact JSON, a separate checker and a formalization
plan are supplied; no Lean build or independent mathematical review is claimed.

\begingroup\small
\begin{thebibliography}{4}
\bibitem{SZ} The Clankers (2026). \emph{Mixed-boundary recovery on the original
SplitZero quotient}. Zeta workbench, \texttt{tau-mixed-boundary-formal/RESEARCH\_NOTE.md},
BR1--BR2 and BR6--BR8; read at \texttt{1efd5333\ldots2644c3}.
The complete commit and source hash are retained in the accompanying ledger.
\bibitem{GH} Greene, J., \& Higgins, C. (2025). Maximal subset sums in a group
of order a power of 2. \emph{Journal of Integer Sequences, 28}, 25.7.8,
\S1 and Theorem 12. General cyclic-cover background, not an assertion that
all integer packets here are perfect covers.
\bibitem{ET} Elsholtz, C., \& Tao, T. (2013). Counting the number of solutions
to the Erd\H{o}s--Straus equation on unit fractions.
\emph{Journal of the Australian Mathematical Society, 94}(1), 50--105.
\S2; arXiv:1107.1010.
\bibitem{Dirichlet} Sutherland, A. V. (2017). \emph{Dirichlet L-functions,
primes in arithmetic progressions}. MIT 18.785, Lecture 18, Theorem 18.1.
\end{thebibliography}
\endgroup
\end{document}
